\chapter{Теоретическая часть. Описание алгоритмов}

\section{Шум Перлина}

Шум Перлина был предложен Кеном Перлином в 1982 году как способ получения плавной псевдослучайной функции, пригодной для компьютерной графики и процедурной генерации \cite{perlin_eevee,perlin_habr,perlin_rtouti}. В отличие от обычного белого шума, соседние значения в шуме Перлина изменяются плавно, поэтому изображение не выглядит набором независимых случайных точек.

В общем виде шум Перлина можно рассматривать как функцию
\[
    Perlin(\overline{x}, seed): \mathbb{R}^n \rightarrow \mathbb{R},
\]
где $\overline{x}$ --- точка пространства, а $seed$ задаёт начальное состояние генератора псевдослучайных значений. На практике чаще всего используются двумерный и трёхмерный случаи, так как они наиболее востребованы при построении текстур, карт высот и анимаций.

Основная идея алгоритма состоит в следующем:
\begin{enumerate}
    \item Пространство разбивается на регулярную сетку из $n$-мерных кубов.
    \item Каждой вершине сетки сопоставляется псевдослучайный градиентный вектор.
    \item Для заданной точки определяется ячейка сетки, в которой она находится.
    \item Для каждой вершины этой ячейки вычисляется вектор от вершины к заданной точке.
    \item Находятся скалярные произведения градиентных векторов и соответствующих векторов расстояния.
    \item Полученные значения интерполируются по координатам точки внутри ячейки.
\end{enumerate}

Для двумерного случая в качестве градиентов могут использоваться, например, следующие векторы:
\[
\{(0,1),(0,-1),(1,0),(-1,0)\}.
\]
Для трёхмерного случая часто используют набор из двенадцати направлений:
\[
\left\{
\begin{aligned}
&(1,1,0),(-1,1,0),(1,-1,0),(-1,-1,0),\\
&(1,0,1),(-1,0,1),(1,0,-1),(-1,0,-1),\\
&(0,1,1),(0,-1,1),(0,1,-1),(0,-1,-1)
\end{aligned}
\right\}.
\]
Такие векторы удобны тем, что имеют одинаковую длину и достаточно равномерно задают направления в пространстве.

Для сглаживания координат внутри ячейки используется функция затухания:
\[
    \varphi(t)=6t^5-15t^4+10t^3.
\]
Она удовлетворяет условиям $\varphi(0)=0$, $\varphi(1)=1$, а её производная в концах отрезка равна нулю. Благодаря этому переходы между соседними ячейками становятся гладкими.

После вычисления скалярных произведений проводится линейная интерполяция. В двумерном случае сначала интерполируются значения по оси $x$, затем результат интерполируется по оси $y$. В трёхмерном случае добавляется аналогичная интерполяция по оси $z$. Итоговое значение можно использовать как яркость пикселя или как высоту точки на карте рельефа.

\section{Симплексный шум}

Симплексный шум был предложен Кеном Перлином как развитие классического шума Перлина \cite{simplex_noise}. Основное отличие состоит в том, что пространство разбивается не на квадраты или кубы, а на симплексы. В двумерном случае симплексом является треугольник, в трёхмерном --- тетраэдр.

Функцию симплексного шума можно записать в виде
\[
    Simplex(\overline{x}, seed): \mathbb{R}^n \rightarrow \mathbb{R}.
\]
Преимущество алгоритма состоит в том, что для вычисления значения в точке нужно учитывать не $2^n$ вершин гиперкуба, как в классическом шуме Перлина, а только $n+1$ вершину симплекса. Поэтому симплексный шум лучше масштабируется при увеличении размерности.

Алгоритм использует преобразование координат, которое называется skew/unskew. Сначала исходное пространство сдвигается так, чтобы регулярная сетка могла быть разбита на симплексы. Для $n$-мерного случая используются коэффициенты
\[
    F_n=\frac{\sqrt{n+1}-1}{n}, \qquad
    G_n=\frac{1-\frac{1}{\sqrt{n+1}}}{n}.
\]
Для точки $\overline{x}=(x_1,x_2,\ldots,x_n)$ вычисляется величина
\[
    s=(x_1+x_2+\ldots+x_n)F_n,
\]
после чего координаты переводятся в скошенное пространство. Затем определяется ячейка, в которой находится точка, и порядок координат внутри ячейки. Этот порядок позволяет определить, в каком именно симплексе находится точка.

После определения вершин симплекса каждой вершине назначается градиентный вектор. Далее для каждой вершины вычисляется вклад:
\[
    contribution = t^4 (\overline{g}\cdot\overline{d}),
\]
где $\overline{g}$ --- градиентный вектор, $\overline{d}$ --- вектор от вершины к точке, а $t$ --- величина, зависящая от расстояния до вершины. Если точка находится слишком далеко от вершины, вклад этой вершины считается равным нулю.

Сумма вкладов всех вершин симплекса образует итоговое значение шума. Симплексный шум часто применяется для генерации плавных текстур, облаков, воды и анимаций, так как хорошо работает не только в двумерном, но и в трёхмерном или четырёхмерном пространстве.

\section{Алгоритм Diamond-Square}

Алгоритм Diamond-Square применяется для генерации фрактальных карт высот. Он был предложен Фурнье, Фусселем и Карпентером и широко используется для процедурного построения рельефа. Алгоритм работает с квадратной матрицей размера
\[
    2^n + 1.
\]
Такой размер необходим для последовательного деления квадратов пополам.

Ход алгоритма:
\begin{enumerate}
    \item Создаётся матрица высот размера $2^n+1$ на $2^n+1$.
    \item В четырёх углах матрицы задаются начальные случайные значения.
    \item Выполняется шаг diamond: для каждого квадрата вычисляется центральная точка как среднее значение четырёх углов плюс случайное смещение.
    \item Выполняется шаг square: для середины каждой стороны вычисляется среднее значение соседних точек плюс случайное смещение.
    \item Размер рассматриваемого квадрата уменьшается в два раза, а амплитуда случайного смещения уменьшается.
    \item Шаги повторяются до тех пор, пока не будут заполнены все точки матрицы.
\end{enumerate}

Параметр roughness определяет силу случайных смещений. Чем больше roughness, тем более резким и гористым получается рельеф. При меньших значениях поверхность становится более гладкой. После завершения алгоритма значения обычно нормируются в диапазон от $0$ до $1$, что удобно для дальнейшей визуализации.

\section{Клеточные автоматы}

Клеточный автомат --- это дискретная модель, состоящая из решётки клеток, каждая из которых находится в одном из конечного числа состояний \cite{cellular_habr,cellular_donntu}. В данной работе рассматриваются двумерные бинарные клеточные автоматы, где каждая клетка может иметь состояние $0$ или $1$. Например, значение $0$ может соответствовать воде или пустоте, а значение $1$ --- суше или стене.

Состояние клетки на следующем шаге зависит от её текущего состояния и состояний соседних клеток. В работе используются два типа окрестностей:
\begin{enumerate}
    \item Окрестность Мура --- все клетки вокруг выбранной клетки в квадрате радиуса $r$. При $r=1$ такая окрестность содержит 8 соседей.
    \item Окрестность фон Неймана --- клетки, расстояние до которых по Манхэттенской метрике не превышает заданный радиус $r$. При $r=1$ такая окрестность содержит 4 соседа: сверху, снизу, слева и справа.
\end{enumerate}

Правило клеточного автомата удобно записывать в виде строки
\[
    B/S,
\]
где $B$ задаёт количества соседей, при которых мёртвая клетка становится живой, а $S$ --- количества соседей, при которых живая клетка сохраняет своё состояние. Например, правило вида
\[
    B.3.6.7.8/S.3.4.6.7.8
\]
означает, что клетка со значением $0$ переходит в состояние $1$ при наличии 3, 6, 7 или 8 соседей, а клетка со значением $1$ остаётся живой при наличии 3, 4, 6, 7 или 8 соседей.

Клеточные автоматы полезны при генерации пещер, островов и других бинарных структур. Начальная матрица может быть случайной, а после нескольких итераций применения правила она приобретает более связную и естественную форму.

\section{Многооктавный шум}

Для получения более сложной и естественной структуры можно использовать многооктавный шум. Он строится как сумма нескольких шумовых слоёв. На каждой следующей октаве частота увеличивается, а амплитуда уменьшается. В данной работе используется схема:
\[
    frequency_{i+1}=2frequency_i, \qquad amplitude_{i+1}=\frac{1}{2}amplitude_i.
\]
Итоговое значение вычисляется как взвешенная сумма:
\[
    noise = \sum_{i=0}^{k-1} amplitude_i \cdot noise_i(frequency_i \cdot x, frequency_i \cdot y),
\]
где $k$ --- количество октав. После суммирования результат нормируется в диапазон от $0$ до $1$.

Многооктавность позволяет совместить крупные формы рельефа и мелкие детали. Такой подход применяется для шума Перлина, симплексного шума и алгоритма Diamond-Square.

\chapter{Практическая часть. Реализация приложения на Python}

\section{Общая структура приложения}

Практическая часть работы представляет собой приложение на языке Python с графическим интерфейсом. Для построения окна используется библиотека Tkinter, а для отображения матриц применяется Matplotlib. Результат работы каждого алгоритма хранится в виде двумерного массива NumPy. Значение элемента массива интерпретируется как яркость пикселя при отображении функцией \texttt{imshow} с цветовой картой \texttt{gray}.

Главное окно приложения состоит из области визуализации и панели управления. В области визуализации отображается текущая матрица. На панели управления находятся кнопки запуска алгоритмов: шум Перлина, симплексный шум, клеточный автомат, Diamond-Square, бинаризация, применение маски и сложение с другим шумом.

Для перерисовки изображения используется функция, которая очищает оси Matplotlib, заново вызывает \texttt{imshow}, скрывает координатные оси и обновляет объект \texttt{canvas}. Такой подход позволяет изменять изображение без перезапуска приложения.

\section{Реализация шума Перлина}

Реализация шума Перлина вынесена в отдельный модуль \texttt{Perlin.py}. Базовая функция принимает координаты точки, значение \texttt{seed} и размер таблицы перестановок. Для каждой точки определяется ячейка сетки, вычисляются локальные координаты, затем значения в вершинах интерполируются с использованием функции сглаживания.

В приложении пользователь задаёт параметры:
\begin{enumerate}
    \item ширину области по оси $x$;
    \item высоту области по оси $y$;
    \item количество пикселей на одну единицу координат;
    \item seed;
    \item размер таблицы;
    \item количество октав.
\end{enumerate}

По введённым параметрам создаются массивы координат с помощью \texttt{numpy.linspace}, затем строится сетка \texttt{meshgrid}. После этого вызывается функция генерации фрактального шума Перлина. Многооктавность реализуется как суммирование нескольких слоёв: каждый следующий слой вычисляется с удвоенной частотой и уменьшенной амплитудой. В конце результат нормируется.

\section{Реализация симплексного шума}

Симплексный шум реализован в модуле \texttt{SimplexNoise.py}. В модуле есть функция генерации нормированных градиентов для $n$-мерного случая. Для вычисления значения шума используются skew/unskew-преобразования, после чего определяется симплекс, в котором находится точка. Далее вычисляются вклады вершин симплекса и суммируются.

В графическом интерфейсе пользователь задаёт seed, размерность пространства от 2 до 4, размеры двумерной плоскости, а также значения координат $z$ и $w$ для среза в трёхмерном или четырёхмерном пространстве. Если выбрана размерность 2, строится обычная двумерная карта шума. Если размерность равна 3 или 4, приложение строит двумерный срез в пространстве большей размерности.

Для двумерного случая используется многооктавная генерация. Это позволяет получить более детализированную текстуру: первая октава формирует крупные области, а последующие добавляют мелкие изменения.

\section{Реализация алгоритма Diamond-Square}

Алгоритм Diamond-Square реализован в модуле \texttt{DiamondSquare.py}. Пользователь задаёт параметр $n$, seed, roughness и количество октав. Размер результирующей карты равен $2^n+1$.

Функция сначала инициализирует четыре угловые точки случайными значениями. Затем в цикле выполняются шаги diamond и square. На каждом уровне детализации размер обрабатываемого квадрата уменьшается в два раза, а амплитуда случайного смещения уменьшается. После заполнения карты выполняется нормализация.

Для многооктавной версии создаётся несколько карт высот. Каждая следующая карта содержит более высокую частоту деталей. Если размеры карт не совпадают, дополнительный слой приводится к базовому размеру. После этого карты складываются с учётом амплитуд и нормируются.

\section{Реализация клеточного автомата}

Клеточный автомат реализован в модуле \texttt{Cellar\_automat.py}. В качестве входных данных используется бинарная матрица. Для каждой клетки вычисляется количество соседей в выбранной окрестности: Мура или фон Неймана. Поддерживается параметр расстояния, который определяет радиус поиска соседей.

Правило автомата задаётся строкой. Строка разделяется на две части: условия рождения и условия сохранения клетки. После подсчёта соседей для каждой клетки формируется новая матрица. Если клетка была равна $0$ и число соседей входит в множество рождения, она становится равной $1$. Если клетка была равна $1$ и число соседей входит в множество сохранения, она остаётся равной $1$. В остальных случаях клетка становится равной $0$.

Такой подход позволяет преобразовывать случайную бинарную матрицу в более связную структуру. Это полезно для генерации пещер, островов и областей суши.

\section{Комбинирование шумов и постобработка}

В приложении реализовано комбинирование нескольких шумовых карт. При добавлении второго шума открывается отдельное окно, в котором можно сгенерировать дополнительную матрицу. Если её размер отличается от размера основной матрицы, она масштабируется до нужной формы. После этого выполняется усреднение:
\[
    matrix = \frac{matrix + matrix_2}{2}.
\]
Такой способ позволяет смешивать различные алгоритмы, например добавить к карте Diamond-Square слой шума Перлина или симплексного шума.

Также реализована маска острова. Она строится по формуле
\[
    z=\frac{1}{(x-x_c)^2+(y-y_c)^2+1},
\]
где $(x_c,y_c)$ --- центр матрицы. При поэлементном умножении на такую маску значения в центре остаются большими, а ближе к краям уменьшаются. Это позволяет получать форму острова.

Дополнительно реализована бинаризация изображения. Пользователь задаёт порог, после чего все значения меньше порога становятся равными $0$, а остальные значения становятся равными $1$:
\[
    a_{ij}=\begin{cases}
    0, & a_{ij}<b,\\
    1, & a_{ij}\ge b.
    \end{cases}
\]
Эта операция необходима перед применением клеточного автомата, так как автомат работает с бинарными состояниями.

\section{Вывод по практической части}

В результате практической части было создано приложение, позволяющее запускать несколько алгоритмов генерации, настраивать их параметры и визуально сравнивать результаты. Использование NumPy позволило представить изображение как матрицу чисел, а Matplotlib обеспечил удобную визуализацию. Tkinter использован для создания окон ввода параметров и управления алгоритмами.

Разработанная программа демонстрирует, что разные методы процедурной генерации дают различные типы структур. Шум Перлина и симплексный шум удобны для плавных текстур, Diamond-Square подходит для карт высот, а клеточные автоматы позволяют получать дискретные формы. Комбинирование этих методов расширяет возможности генерации и позволяет создавать более сложные изображения.
